\documentclass{amsart}
\usepackage{amsmath,amssymb}
\usepackage[hidelinks]{hyperref}
\hypersetup{
  pdftitle={Equal radicals of consecutive products},
  pdfauthor={Ruiyi Zhang}
}

\newtheorem{theorem}{Theorem}
\newtheorem{lemma}{Lemma}[section]
\newtheorem{proposition}[lemma]{Proposition}
\newtheorem{corollary}[lemma]{Corollary}
\theoremstyle{remark}
\newtheorem{remark}[lemma]{Remark}
\numberwithin{equation}{section}

\title{Equal radicals of consecutive products}
\author{Ruiyi Zhang}
\address{Independent Researcher, Columbia, Maryland, USA}
\email{ryzfsu@gmail.com}
\subjclass[2020]{Primary 11A25; Secondary 11D25}
\keywords{radicals, consecutive integers, Diophantine equations}

\begin{document}

\begin{abstract}
For every fixed $k\ge2$ and $\varepsilon>0$,
$F_{k,k}(X)\ll_{k,\varepsilon}X^{2/3+\varepsilon}$. For $k\ge3$ this
improves the $X^{1+o(1)}$ bound of Lebowitz--Lockard. For $k=2$, Kane's
localized estimates give the large-quotient endpoint bound. Our independent
proof combines a Pell bound for partner multiplicities, a quotient-tail
estimate, and Farey spacing for small quotients. Reuss's determinant estimate
also gives
$N(H,Y)\ll_\varepsilon\allowbreak H^\varepsilon(Y/H+H^{7/12})$, uniformly for
$H\ge2$ and $1\le Y<H^2$, where $N(H,Y)$ counts integers
$t\in(H,2H]$ with $\operatorname{rad}(t(t+1))\le Y$. A cubic family gives
$F(X)\ge(4^{1/3}/3)X^{1/3}-O(1)$.
\end{abstract}

\maketitle

\section{Introduction}\label{sec:introduction}

For a positive integer $a$, let $\operatorname{rad}(a)=\prod_{p\mid a}p$.
We study the ordered count
\[
 F(X)=\#\{1\le m<n\le X:
 \operatorname{rad}(m(m+1))=\operatorname{rad}(n(n+1))\}.
\]
More generally, put
\[
 P_k(t)=\prod_{j=0}^{k-1}(t+j),\qquad
 F_{k,k}(X)=\#\{1\le m<n\le X:
 \operatorname{rad}(P_k(m))=\operatorname{rad}(P_k(n))\}.
\]
Thus $F(X)=F_{2,2}(X)$. We impose no disjointness condition on the two
blocks, and all implied constants may depend on the fixed integer $k$.

\begin{theorem}\label{thm:main}
For every fixed integer $k\ge2$, every $\varepsilon>0$, and every real
$X\ge2$,
\[
 F_{k,k}(X)\ll_{k,\varepsilon}X^{2/3+\varepsilon}.
\]
\end{theorem}

Erd\H{o}s asked for the size of $F(X)$ \cite{Erdos1980}.
Balasubramanian, Shorey and Waldschmidt \cite[\S5]{BSW1989} studied
non-overlapping blocks of consecutive integers with the same prime divisors.
Lebowitz--Lockard \cite[Theorems~5 and~10]{LL2025} proved
$F_{k,k}(X)\le X^{1+o(1)}$. For $k\ge3$, Theorem~\ref{thm:main}
replaces this exponent by $2/3+\varepsilon$. Pell equations of this kind
go back to Lehmer's work on St{\o}rmer's problem \cite{Lehmer1964}; see
\cite[\S2]{BSW1989}.

For $k=2$, Kane's localized estimates \cite[Propositions~7 and~10]{Kane2014}
(Propositions~6 and~9 of \cite{Kane2011}) give the large-quotient endpoint
bound after specialization to $(1,t,-(t+1))$ and dyadic summation. In the
proof below, Lemma~\ref{lem:small-quotient} handles the small-quotient range;
together with the endpoint-fibre and quotient-tail bounds, it gives an
independent proof of the $k=2$ case.

Put $q_k(n)=P_k(n)/\operatorname{rad}(P_k(n))$. Splitting at a cutoff $K$
and applying the endpoint-fibre, quotient-tail and small-quotient estimates
gives, for $\max\{k,2\}\le K\le X$,
\begin{equation}\label{eq:introduction-split}
F_{k,k}(X)\ll_{k,\varepsilon}
X^\varepsilon\bigl(K^2+X/K+X^{1/2}\bigr).
\end{equation}
For $X\ge k^3$, taking $K=X^{1/3}$ in
\eqref{eq:introduction-split} proves Theorem~\ref{thm:main}. The bounded
range $2\le X<k^3$ is absorbed into the implied constant, since $k$ is fixed.
Theorem~\ref{thm:cubic-family} gives an explicit lower-bound family for
Erd\H{o}s's count. Theorem~\ref{thm:low-radical} follows from known estimates;
its deduction appears in Appendix~A.
Section~\ref{sec:consequences} gives two consequences of the pair bound.
The tools used for numerical checks and literature searches are identified
at the end of the paper.

\section{Proof of the fixed-block theorem}\label{sec:fixed-block}

Fix an integer $k\ge2$. Put
\[
P_k(n)=\prod_{j=0}^{k-1}(n+j),\qquad
q_k(n)=\frac{P_k(n)}{\operatorname{rad}(P_k(n))}.
\]
The function $P_k$ is positive and strictly increasing on the positive
integers. Write $\omega(D)$ for the number of distinct prime divisors of
$D$. All implied constants in this section may depend on $k$.

Lemma~2.1 is a Pell-equation count; compare Lebowitz--Lockard
\cite[Lemma~8]{LL2025} and Tao \cite[Lemma~2.10]{Tao2026}.

\begin{lemma}\label{lem:support-fibre}
For every squarefree positive integer $Q$ and every real $X\ge2$,
\[
\#\{1\le z\le X:\operatorname{rad}(z(z+1))\mid Q\}
\ll 2^{\omega(Q)}\log(2X),
\]
with an absolute implied constant. Consequently, for every $\delta>0$,
every value of $\operatorname{rad}(P_k(n))$ has at most
\begin{equation}\label{eq:endpoint-fibre}
C_{k,\delta}X^{k\delta}\log(2X)
\end{equation}
preimages among the positive integers $n\le X$.
\end{lemma}

\begin{proof}
Write $z(z+1)=dy^2$, where $d$ is squarefree. Then $d\mid Q$ and
\[
(2z+1)^2-d(2y)^2=1.
\]
The case $d=1$ has no solution with $z\ge1$, since
$(2z+1-2y)(2z+1+2y)=1$ is impossible. For a fixed $d\ge2$, put
$A=2z+1$, $B=2y$ and $\alpha=A+B\sqrt d$. Thus
\[
1<\alpha<2A\le4X+2,\qquad \alpha^{-1}=A-B\sqrt d.
\]
If $\alpha_2>\alpha_1$ arise from two solutions, then
\[
\frac{\alpha_2}{\alpha_1}=U_0+V_0\sqrt d>1,
\qquad U_0,V_0\in\mathbb Z,\qquad U_0^2-dV_0^2=1.
\]
Its conjugate is its positive reciprocal. Hence
\[
U_0=\frac12\left(\frac{\alpha_2}{\alpha_1}
                 +\frac{\alpha_1}{\alpha_2}\right)\ge2,
\qquad
\frac{\alpha_2}{\alpha_1}\ge2+\sqrt3.
\]
There are $O(\log(2X))$ such solutions for each $d\mid Q$. Summing over
the squarefree divisors of $Q$ proves the first assertion.

For every $\delta>0$,
\[
2^{\omega(Q)}\ll_\delta Q^\delta.
\]
Indeed, $2\le p^\delta$ for all sufficiently large primes $p$, and the
remaining primes contribute a constant. If
$Q=\operatorname{rad}(P_k(n))$ with $n\le X$, then
$Q\le(X+k-1)^k\ll_k X^k$. Every preimage $m$ of $Q$ satisfies
$\operatorname{rad}(m(m+1))\mid Q$, since $k\ge2$. This gives
\eqref{eq:endpoint-fibre}.
\end{proof}

Erd\H{o}s and Straus \cite[p.~63]{ErdosStraus1977} also use a Pell equation
for products of consecutive integers.

For a positive integer $D$, let
\[
\rho_k(D)=\#\{x\bmod D:P_k(x)\equiv0\pmod D\}.
\]

\begin{lemma}\label{lem:block-roots}
For every positive integer $D$,
\[
\rho_k(D)\le(k-1)!\,k^{\omega(D)}.
\]
\end{lemma}

\begin{proof}
If $p>k-1$, the $k$ factors of $P_k$ have distinct roots modulo $p$.
At most one factor is divisible by $p$, so $\rho_k(p^j)=k$ for every
$j\ge1$.

For $p\le k-1$, put $C_p=v_p((k-1)!)$. Given an integer $x$, choose
$i\in\{0,\ldots,k-1\}$ for which $v_p(x+i)$ is maximal, taking
$v_p(0)=\infty$. For $j\ne i$, the ultrametric inequality and maximality
give
\[
v_p(x+j)\le v_p(j-i).
\]
Thus
\[
\sum_{j\ne i}v_p(x+j)
\le v_p\bigl(i!(k-1-i)!\bigr)\le C_p.
\]
If $p^j\mid P_k(x)$, then
$p^{\max(j-C_p,0)}\mid x+i$. For each $i$ there are at most $p^{C_p}$
such residue classes modulo $p^j$. Hence
$\rho_k(p^j)\le kp^{C_p}$. The Chinese remainder theorem now gives
\[
\rho_k(D)\le k^{\omega(D)}\prod_{p\le k-1}p^{C_p}
=(k-1)!\,k^{\omega(D)}.
\qedhere
\]
\end{proof}

Sanna \cite[Lemma~3.3]{Sanna2022} proves the related bound
$O_r(x/y^{1/2})$ for the number of $n\le x$ such that one of
$n,\ldots,n+r-1$ has powerful part exceeding $y$.

\begin{lemma}\label{lem:quotient-tail}
For all real $X\ge2$, $K\ge\max(k,2)$ and $0<\eta<1$,
\begin{equation}\label{eq:quotient-tail}
\#\{1\le n\le X:q_k(n)>K\}
\ll_{k,\eta}XK^{-1+\eta}
       +K^2(1+\log K)^{k-1}+X^{1/2}.
\end{equation}
\end{lemma}

\begin{proof}
Let $1\le n\le X$ satisfy $q_k(n)>K$. Suppose first that a prime $p>K$
divides $q_k(n)$. Then
$p^2\mid P_k(n)$. Since $p>K\ge k$, exactly one factor is divisible by
$p$, and $p^2\mid n+i$ for some $0\le i<k$. The number of these $n$ is
at most
\[
k\sum_{K<p\le\sqrt{X+k-1}}\left(\frac X{p^2}+1\right)
\ll_k\frac XK+X^{1/2}.
\]
The last bound follows by summing over all integers in place of primes.

Otherwise every prime factor of $q_k(n)$ is at most $K$. Multiply its
prime factors, with multiplicity, until the product first exceeds $K$.
This gives an integer $g$ such that
\[
g\mid q_k(n),\qquad K<g\le K^2,
\qquad g\operatorname{rad}(g)\mid P_k(n).
\]
For the last divisibility, if $p^a\mid g$, then
$a\le v_p(P_k(n))-1$. Lemma~\ref{lem:block-roots} bounds the number of
these $n$ by
\begin{equation}\label{eq:tail-crt}
(k-1)!\sum_{K<g\le K^2}k^{\omega(g)}
\left(\frac{X}{g\operatorname{rad}(g)}+1\right).
\end{equation}
For $0<\eta<1$, the Euler product
\[
\sum_{g\ge1}\frac{k^{\omega(g)}g^{1-\eta}}
                   {g\operatorname{rad}(g)}
=\prod_p\left(1+\frac{k}{p(p^\eta-1)}\right)
\]
converges, since its local factor is $1+O_{k,\eta}(p^{-1-\eta})$.
The first term of \eqref{eq:tail-crt} is therefore
$O_{k,\eta}(XK^{-1+\eta})$.

Let $\tau_k(g)$ count the ordered factorizations of $g$ into $k$
positive integers. Then $k^{\omega(g)}\le\tau_k(g)$ and, for $Y\ge1$,
\[
\sum_{g\le Y}\tau_k(g)
\le Y\left(\sum_{r\le Y}\frac1r\right)^{k-1}
\le Y(1+\log Y)^{k-1}.
\]
The second term of \eqref{eq:tail-crt} is
$O_k(K^2(1+\log K)^{k-1})$. This proves \eqref{eq:quotient-tail}.
\end{proof}

\begin{lemma}\label{lem:small-quotient}
For every real $X\ge2$ and $K\ge1$,
\[
\#\left\{(m,n):
\begin{gathered}
1\le m<n\le X,\quad q_k(m),q_k(n)\le K,\\
\operatorname{rad}(P_k(m))=\operatorname{rad}(P_k(n))
\end{gathered}\right\}
\ll_k K^2\log(2X).
\]
\end{lemma}

\begin{proof}
Put $\mu=q_k(m)$ and $\nu=q_k(n)$. Strict monotonicity of $P_k$ gives
\begin{equation}\label{eq:quotient-curve}
1\le\mu<\nu\le K,\qquad \nu P_k(m)=\mu P_k(n).
\end{equation}
Set
\[
\xi=2m+k-1,\qquad \zeta=2n+k-1,
\qquad
R_k(U)=\prod_{j=0}^{k-1}(U+2j-k+1).
\]
Then $R_k(\zeta)=2^kP_k(n)$. With
\[
\mathcal H_k=\{h:1\le h\le k-1,\ h\equiv k-1\pmod2\},
\qquad A_k=\sum_{h\in\mathcal H_k}h^2,
\]
we have
\[
R_k(U)=U^kG_k(U),\qquad
G_k(U)=\prod_{h\in\mathcal H_k}\left(1-\frac{h^2}{U^2}\right).
\]
For $U\ge k+1$, the function $G_k$ is strictly increasing and
\begin{equation}\label{eq:centered-defect}
0<G_k(U)<1,\qquad 1-G_k(U)\le\frac{A_k}{U^2}.
\end{equation}
Indeed, $\mathcal H_k$ is nonempty and each factor is positive and
strictly increasing. The last inequality follows from
$1-\prod_j(1-x_j)\le\sum_jx_j$ for $0\le x_j\le1$.

Fix $\mu,\nu$ in \eqref{eq:quotient-curve}, and put
$\gamma=(\nu/\mu)^{1/k}$ and $\lambda=\zeta/\xi$.
Thus $\zeta>\xi\ge k+1$ and
\[
\gamma^k=\lambda^k\frac{G_k(\zeta)}{G_k(\xi)},\qquad 1<\lambda<\gamma.
\]
Using \eqref{eq:centered-defect}, we obtain
\begin{align*}
0<\gamma-\lambda
&=\gamma\left(1-\left(\frac{G_k(\xi)}{G_k(\zeta)}\right)^{1/k}\right)\\
&\le\gamma\bigl(1-G_k(\xi)^{1/k}\bigr)\\
&\le\gamma\bigl(1-G_k(\xi)\bigr)\\
&\le\frac{A_k\gamma}{\xi^2}.
\end{align*}
Write $\lambda=h/q$ in lowest terms, with $h,q>0$. Since $q\mid\xi$,
\[
0<\gamma-\frac hq\le\frac{A_k\gamma}{q^2},\qquad q\le2X+k-1.
\]
In a block $2^j\le q<2^{j+1}$, these fractions lie in an interval of
length $A_k\gamma/2^{2j}$. Distinct reduced fractions in the block are
more than $2^{-2j-2}$ apart. There are at most $1+4A_k\gamma$ fractions
in each block, and hence $O_k(\gamma\log(2X))$ possible values of $\lambda$.

For each such $\lambda$, every possible $\xi$ is a root of
\[
\nu R_k(U)-\mu R_k(\lambda U)=0.
\]
This polynomial has degree $k$, since its leading coefficient is
$\nu-\mu\lambda^k=\mu(\gamma^k-\lambda^k)>0$.
It has at most $k$ real roots. Each root determines
at most one pair $(m,n)$. Therefore, for fixed $\mu,\nu$, the number of
pairs in \eqref{eq:quotient-curve} is
\[
\ll_k (\nu/\mu)^{1/k}\log(2X).
\]
Finally,
\[
\sum_{1\le\mu<\nu\le K}(\nu/\mu)^{1/k}
\le\sum_{\nu\le K}\nu^{1/k}\frac{k}{k-1}\nu^{1-1/k}
\ll_k K^2.
\qedhere
\]
\end{proof}

\begin{proof}[Proof of Theorem~\ref{thm:main}]
Let $\varepsilon>0$. A counted pair either has both quotients at most
$K$, or has an endpoint in the exceptional set of
Lemma~\ref{lem:quotient-tail}. Multiplying the latter count by
\eqref{eq:endpoint-fibre} and applying Lemma~\ref{lem:small-quotient}
gives, for $K\ge\max(k,2)$,
\begin{align}\label{eq:fixed-block-assembly}
F_{k,k}(X)\ll_{k,\delta,\eta}{}&K^2\log(2X)\notag\\
&+X^{k\delta}\log(2X)
\left(XK^{-1+\eta}+K^2(1+\log K)^{k-1}+X^{1/2}\right).
\end{align}
The fibre bound counts all possible partners $1\le m\le X$, so this
also covers pairs whose smaller endpoint is exceptional.

Ignoring the factors $X^{k\delta}$, $K^\eta$ and the logarithms,
$K=X^\tau$ gives the exponent $\max(2\tau,1-\tau,1/2)$, minimized at
$\tau=1/3$. Set $K=\max(k,2,X^{1/3})$. If $X\ge k^3$, then
$K=X^{1/3}$. With this
choice fixed, take
\[
\eta=\min\left(\frac12,\frac\varepsilon2\right),
\qquad \delta=\frac\varepsilon{4k}.
\]
Every power of $X$ in \eqref{eq:fixed-block-assembly} is at most
\[
\frac23+k\delta+\frac\eta3
\le\frac23+\frac{5\varepsilon}{12}.
\]
The logarithmic factors together are
$O_{k,\varepsilon}(X^{\varepsilon/2})$. Hence
\[
F_{k,k}(X)\ll_{k,\varepsilon}X^{2/3+\varepsilon}.
\]
The bounded range $2\le X<k^3$ is absorbed into the implied constant.
\end{proof}

\section{A cubic lower-bound family}\label{sec:cubic-family}

\begin{theorem}\label{thm:cubic-family}
There is an absolute constant $C$ such that, for every real $X\ge2$,
\[
F(X)\ge \frac{4^{1/3}}{3}X^{1/3}-C.
\]
\end{theorem}

\begin{proof}
Erd\H{o}s posed the product-support counting problem in
\cite{Erdos1980}. We know of no published source for this particular
polynomial family.
Take integers $t\ge7$ satisfying
\[
t\not\equiv0\pmod4,\qquad t\equiv1,2,7,8\pmod9,
\]
and put
\[
m=\frac{t^2-4}{3},\qquad n=\frac{t^3-3t-2}{4}.
\]
The resulting integers are positive, and
\[
12(n-m)=3t^3-4t^2-9t+10>0
\]
for $t\ge7$, so $m<n$. Direct multiplication gives
\[
m(m+1)=\frac{(t^2-4)(t^2-1)}9,\qquad
n(n+1)=\frac{(t^2-4)(t^2-1)^2}{16}.
\]
For every prime other than $2$ and $3$, the two products have the same
prime support. If $t\equiv\pm1\pmod9$, then $9\mid t^2-1$ and
$3\mid t^2-4$; if $t\equiv\pm2\pmod9$, then $9\mid t^2-4$ and
$3\mid t^2-1$. Thus division by $9$ leaves a factor $3$ in the first
product, and the second product is also divisible by $3$.

If $t$ is odd, then $8\mid t^2-1$, so both products remain even after
their displayed divisions. If $t\equiv2\pmod4$, then
$32\mid t^2-4$, which again leaves both products even after division by
$16$. Hence the two products have equal radicals.

The conditions specify twelve residue classes modulo $36$, of density
$1/3$. Since $n=t^3/4+O(t)$, the number of such $t$ with $n\le X$ is
$(4^{1/3}/3)X^{1/3}+O(1)$. The map $t\mapsto m$ is strictly increasing
for $t\ge7$, so these give distinct pairs counted by $F(X)$.
\end{proof}

\section{A radical-threshold estimate from known results}\label{sec:low-radical}

For real $H\ge2$ and $1\le Y<H^2$, define
\[
N(H,Y)=\#\{t\in\mathbb Z:H<t\le2H,
                 \operatorname{rad}(t(t+1))\le Y\}.
\]

\begin{theorem}\label{thm:low-radical}
For every $\varepsilon>0$,
\[
N(H,Y)\ll_\varepsilon H^\varepsilon
       \left(\frac YH+H^{7/12}\right)
\]
uniformly for real $H\ge2$ and $1\le Y<H^2$.
\end{theorem}

Theorem~3 is a consequence of known estimates, stated here in this notation.
In the proof of Kane's Proposition~7 \cite{Kane2014}, the dyadic
lattice count for triples of height $N_0$ is bounded by
\[
N_0^\varepsilon\left[
\frac{\operatorname{rad}(A)\operatorname{rad}(B)\operatorname{rad}(C)}{|C|}
+v(A)v(B)v(C)\right].
\]
For $H<t\le2H$, take $N_0=\lceil2H+1\rceil$ and
$(A,B,C)=(1,t,-(t+1))$. Put
 \[
S_2(n)=\prod_{p^j\parallel n}p^{\lfloor j/2\rfloor},\qquad
T_2(n)=\frac{n}{S_2(n)^2},\qquad
\Pi=S_2(t)S_2(t+1).
\]
Here $v(n)=\prod_{p^2\mid n}p=\operatorname{rad}(S_2(n))$ in Kane's
notation. When $\Pi\le H^{2/3}$, the displayed lattice
inequality gives $H^\varepsilon(Y/H+\Pi)$. For $\Pi>H^{2/3}$,
Kane's Proposition~10 \cite[Proposition~10]{Kane2014} gives the
square-part estimate
\[
H^\varepsilon\left(T_2(t)T_2(t+1)
 +\sqrt{\Pi T_2(t)T_2(t+1)}\right)
\ll H^\varepsilon\left(\frac{H^2}{\Pi^2}+\frac{H}{\sqrt{\Pi}}\right),
\]
since $t(t+1)\ll H^2$. Both ranges are
$O_\varepsilon(H^\varepsilon(Y/H+H^{2/3}))$, with the split point
$\Pi=H^{2/3}$ balancing the terms in Kane's proof. Diao's Lemma~2.1
\cite{Diao2026} gives the $7/12$ estimate in the large-product range.
Appendix~A gives a detailed deduction via Reuss's theorem. Theorem~3 records
this known consequence in the notation used here.

\section{Lebowitz--Lockard's Conjecture 11}

Lebowitz--Lockard state the following conjecture \cite[Conjecture~11]{LL2025}:
\begin{quote}
If $k,\ell>1$ and $(k,\ell)\ne(2,2)$, then the equation
\[
\operatorname{rad}(m(m+1)\cdots(m+k-1))
=\operatorname{rad}(n(n+1)\cdots(n+\ell-1))
\]
only has finitely many solutions.
\end{quote}
Their preceding definition of $F_{k,\ell}$ imposes $m<n$. The conjecture
also fails under this order condition, for $(k,\ell)=(3,2)$ and $(4,2)$.
Equivalent forms of the identities below appear in Erd\H{o}s--Straus
\cite[p.~64]{ErdosStraus1977}. For every positive integer $t$,
\begin{align*}
P_2(t(t+2))&=t(t+2)(t+1)^2,\\
P_2\!\left(\frac{t(t+3)}2\right)&=\frac14 P_4(t).
\end{align*}
Since $t+1\mid P_3(t)$, the first identity shows that
$(m,n)=(t,t(t+2))$ gives equal radicals for $(k,\ell)=(3,2)$. For the
second identity, four consecutive integers have product divisible by
$24$: one is divisible by $4$, a distinct one is even, and one is divisible
by $3$. Thus $P_4(t)/4$ remains divisible by $6$, so division by $4$ leaves
its prime support unchanged.
The integer $t(t+3)/2$ is integral, and the second identity gives equal
radicals for $(m,n)=(t,t(t+3)/2)$ and $(k,\ell)=(4,2)$. Both families have
$m<n$ for every positive integer $t$. Swapping the variables gives
infinite families for $(2,3)$ and $(2,4)$ if no order condition is imposed.
Conjecture~11 is false as literally stated; its intended form is unclear.
Erd\H{o}s and Straus \cite[p.~64]{ErdosStraus1977} conjectured that
$A(n,k)=\rho A(m,\ell)$ has only finitely many solutions with $m\ge n+k$
for each rational $\rho$ when $\min(k,\ell)>1$ and $\max(k,\ell)>2$,
where $A(n,k)=(n+k)!/n!$. The second identity reads
$A(t-1,4)=4A(t(t+3)/2-1,2)$, and
$t(t+3)/2-1>(t-1)+4$ for every integer $t\ge3$. Their conjecture is
therefore also false as literally stated. They print the underlying
four-factor identity on the same page in their divisibility argument.


\section{Further consequences}\label{sec:consequences}

\subsection{Radical fibres and moments}
\begin{remark}
Lemma~\ref{lem:support-fibre} is a growth form of the Pell argument; compare
Lehmer \cite[Theorem~1]{Lehmer1964} and Lebowitz--Lockard
\cite[Lemma~8]{LL2025}. In particular, if
\[
\nu_R(X)=\#\{1\le n\le X:\operatorname{rad}(n(n+1))=R\},
\]
then Lemma~\ref{lem:support-fibre} gives
\[
\nu_R(X)\ll 2^{\omega(R)}\log(2X)\ll_\varepsilon X^\varepsilon
\qquad(R\ll X^2),
\]
so $\max_R\nu_R(X)\ll_\varepsilon X^\varepsilon$. Since
\[
\sum_R\nu_R(X)=\lfloor X\rfloor,\qquad
\sum_R\nu_R(X)(\nu_R(X)-1)=2F(X),
\]
the main theorem gives
\[
\begin{aligned}
\#\{\operatorname{rad}(n(n+1)):1\le n\le X\}
&=\lfloor X\rfloor-O_\varepsilon(X^{2/3+\varepsilon}),\\
\sum_R\nu_R(X)^r
&=\lfloor X\rfloor+O_{r,\varepsilon}(X^{2/3+\varepsilon})
\end{aligned}
\]
for each fixed integer $r\ge2$. For $r\ge3$, summing the following
pointwise bound over $R$ gives
\[
0\le \nu_R(X)^r-\nu_R(X)
\le (r-1)\bigl(\max_S\nu_S(X)\bigr)^{r-2}
       \nu_R(X)(\nu_R(X)-1),
\]
and hence the stated moment estimate.
The case $r=2$ follows from the pair identity above. That identity also
gives
\[
0\le \lfloor X\rfloor-\#\{R:\nu_R(X)>0\}
=\sum_{\nu_R(X)>0}(\nu_R(X)-1)\le F(X).
\]
\end{remark}

\subsection{Short gaps}
\begin{remark}
For $1\le\Delta\le X$, let $G(X,\Delta)$ count the pairs in $F(X)$ with
$n-m\le\Delta$. For a fixed gap $d=n-m\ge2$, every prime dividing the
common radical divides one factor from each of $m(m+1)$ and $n(n+1)$.
Subtracting these factors shows that it divides one of $d-1,d,d+1$.
There are at most $2^{\omega(d(d-1)(d+1))}$ possible common radicals.
Since $d(d-1)(d+1)\le X^3$, take $\delta=\varepsilon/6$ in
$2^{\omega(q)}\ll_\delta q^\delta$. This contributes at most
$O_\varepsilon(X^{\varepsilon/2})$ radicals. In the fibre estimate, take
its exponent parameter to be $\varepsilon/8$; then
$X^{\varepsilon/4}\log(2X)\ll_\varepsilon X^{\varepsilon/2}$. Thus each
gap contributes $O_\varepsilon(X^\varepsilon)$ pairs. Summing over
$d\le\Delta$ gives
\[
G(X,\Delta)\ll_\varepsilon \Delta X^\varepsilon.
\]
For $d=1$, equality of radicals reduces to
$\operatorname{rad}(m)=\operatorname{rad}(m+2)$, which forces $m=2$ and
gives the single pair $(m,n)=(2,3)$: their common radical divides
$\operatorname{rad}(\gcd(m,m+2))\mid2$. It cannot be $1$, so both terms
are powers of $2$, which leaves $m=2$.
\end{remark}

\appendix
\section{Proof of Theorem~3}\label{app:low-radical-proof}

\begin{proof}[Proof of Theorem~\ref{thm:low-radical}]
Write uniquely
\[
t=d^2r,\qquad t+1=f^2s,
\]
where $r,s$ are squarefree positive integers. Then
\begin{equation}\label{eq:low-radical-pell}
f^2s-d^2r=1,\qquad (d,f)=(r,s)=1.
\end{equation}
Put
\[
\alpha_0=\gcd(r,\operatorname{rad}(d)),\qquad
\beta_0=\gcd(s,\operatorname{rad}(f)).
\]
The radical identity is
\begin{equation}\label{eq:low-radical-identity}
\operatorname{rad}(t(t+1))
=\frac{rs\operatorname{rad}(d)\operatorname{rad}(f)}{\alpha_0\beta_0}.
\end{equation}
Since $rs>H^2/(d^2f^2)$, every point counted by $N(H,Y)$ satisfies
\begin{equation}\label{eq:low-radical-modulus}
d^2f^2\alpha_0\beta_0\ge
\frac{H^2\operatorname{rad}(d)\operatorname{rad}(f)}{Y}.
\end{equation}

Set $H_0=H^{7/12}$ and first count the points with $df\le H_0$. For fixed
$d,f,\alpha_0,\beta_0$, the divisibilities $\alpha_0\mid r$ and $\beta_0\mid s$ give
\[
t\equiv0\pmod{d^2\alpha_0},\qquad t\equiv-1\pmod{f^2\beta_0}.
\]
The moduli are coprime, since every prime dividing the first divides $t$
and every prime dividing the second divides $t+1$. Thus the Chinese
remainder theorem gives at most one class modulo $d^2f^2\alpha_0\beta_0$,
containing at most $H/(d^2f^2\alpha_0\beta_0)+1$ integers in the interval. By
\eqref{eq:low-radical-modulus},
\begin{equation}\label{eq:low-radical-small-product}
N_{\mathrm{lo}}\ll
\frac YH\sum_{df\le H_0}
 \frac{\tau(\operatorname{rad}(d))\tau(\operatorname{rad}(f))}
      {\operatorname{rad}(d)\operatorname{rad}(f)}
 +\sum_{df\le H_0}\tau(d)\tau(f).
\end{equation}
For every $\eta>0$, Rankin's trick and the convergent product
\[
\sum_{n\ge1}\frac{\tau(\operatorname{rad}(n))}
                   {\operatorname{rad}(n)n^\eta}
=\prod_p\left(1+\frac{2}{p(p^\eta-1)}\right)<\infty
\]
give
\[
\sum_{df\le H_0}
 \frac{\tau(\operatorname{rad}(d))\tau(\operatorname{rad}(f))}
      {\operatorname{rad}(d)\operatorname{rad}(f)}
\ll_\eta H_0^\eta.
\]
Also $\sum_{df\le H_0}\tau(d)\tau(f)\ll H_0(\log(2H_0))^3$, by summing
the four-fold divisor function. Taking $\eta$ sufficiently small in
terms of $\varepsilon$ and absorbing logarithms gives
\begin{equation}\label{eq:low-radical-low-bound}
N_{\mathrm{lo}}\ll_\varepsilon
H^\varepsilon\left(\frac YH+H^{7/12}\right).
\end{equation}

It remains to count solutions of \eqref{eq:low-radical-pell} with
$df>H_0$; here we discard the radical restriction. Since $r,s\ge1$,
$t\le2H$, and $t+1\le3H$, we have
\begin{equation}\label{eq:low-radical-squareparts}
d\le\sqrt{2H},\qquad f\le\sqrt{3H}.
\end{equation}
Partition the $d,f$ ranges into dyadic boxes and refine each into
multiplicative cells of ratio $\lambda=2^{1/8}$. On a cell write
\[
 S<t\le\lambda S,\qquad \mathcal D<d\le\lambda \mathcal D,\qquad
\mathcal E<f\le\lambda \mathcal E,
\]
and put $x=S/\lambda^2$, $U=x/\mathcal D^2$, and $V=x/\mathcal E^2$. Then
$r=t/d^2\in(U,\lambda^3U)$ and
$s=(t+1)/f^2\in(V,\lambda^3V+\mathcal E^{-2})$. Since
$\lambda^3+\lambda^2/S<2$ for $S\ge2$, we have
\[
U<r<2U,\qquad V<s<2V,
\]
and the variables lie in dyadic boxes with the exact scale identities
\[
U\mathcal D^2=V\mathcal E^2=x\asymp H.
\]
Only a bounded number of such cells is used in each dyadic box. In every
cell,
\begin{equation}\label{eq:low-radical-box-products}
UV\asymp\frac{H^2}{(\mathcal D\mathcal E)^2}.
\end{equation}

Suppose first that $UV\le H^{1/2}$. Fix squarefree $r,s$. If $rs=1$,
\eqref{eq:low-radical-pell} would give $f^2-d^2=1$, which is impossible
for positive $d,f$. Otherwise
\[
(sf)^2-rsd^2=s.
\]
The element
\[
\xi=sf+d\sqrt{rs}
      =\sqrt s\bigl(\sqrt{t+1}+\sqrt t\bigr)
\]
is totally positive and has norm $s$. The relevant quadratic order is
maximal unless $rs\equiv1\pmod4$; in that case it has conductor $2$,
and $s$ is odd and coprime to the conductor. Thus the ideals of norm $s$
in the order are invertible and correspond to ideals of norm $s$ in the
maximal order. Every prime dividing $s$ ramifies, so at most $\tau(s)$ such
ideals arise. Generators of the same principal ideal differ by a unit of
the order; for two totally positive generators their ratio is a totally
positive norm-one unit. If
$\epsilon>1$ is such a unit, its integer trace is at least $3$, whence
$\epsilon\ge(3+\sqrt5)/2>2$. In contrast, the displayed formula for
$\xi$ shows that its ratio across $H<t\le2H$ is less than $\sqrt2$.
Thus each ideal contributes at most one solution. Since $s\le3H$,
$\tau(s)\ll_\varepsilon H^\varepsilon$. There are $O(UV)$ choices of
$r,s$ in their boxes, and therefore
\begin{equation}\label{eq:low-radical-pell-range}
N_{UV\le H^{1/2}}\ll_\varepsilon H^\varepsilon UV
                       \ll_\varepsilon H^{1/2+\varepsilon}.
\end{equation}

Now suppose $UV>H^{1/2}$. Equations
\eqref{eq:low-radical-box-products} and $df>H_0$ give
\[
H^{7/12}\ll \mathcal D\mathcal E\ll H^{3/4},\qquad
\log(\mathcal D\mathcal E)\asymp\log(UV)\asymp\log H.
\]
Apply Reuss's determinant estimate \cite[Theorem~1]{Reuss2014} to
$sf^2-rd^2=1$ with $(k,\ell,h)=(2,1,1)$ and scales
$U=x/\mathcal D^2$, $V=x/\mathcal E^2$. They satisfy
$U\mathcal D^2=V\mathcal E^2=x$. Since $df>H_0$ and
\eqref{eq:low-radical-squareparts} gives $d\le\sqrt{2H}$ and
$f\le\sqrt{3H}$, we have
\[
d>H_0/\sqrt{3H}\gg H^{1/12},\qquad
f>H_0/\sqrt{2H}\gg H^{1/12}.
\]
As $d\le\lambda \mathcal D$, $f\le\lambda \mathcal E$, and $x\asymp H$, this implies
$\mathcal D,\mathcal E,x>1$ for large $H$. Define
\[
M=\exp\!\left(\frac{9\log(\mathcal D\mathcal E)\log(UV)}{8\log x}\right).
\]
In this range $\mathcal D\mathcal E>1$, $UV>H^{1/2}>1$, and $x>1$, so $M>1$.
If the sides are interchanged, the equation has $h=-1$. Reuss's
hypotheses are $1\le\ell<k$, $h\ne0$,
$\log(\mathcal D\mathcal E)\asymp\log(UV)\asymp\log x$, and, when $\ell=1$,
$\mathcal D\mathcal E\gg x^{1/k}$. We have $\mathcal D\mathcal E\gg H^{7/12}\gg x^{1/2}$, and the logarithmic
comparabilities were established above. The estimate applies for large
$x$; bounded values of $H$ are absorbed in the implied constant. The
resulting bound for the number $N$ in a cell is
\begin{equation}\label{eq:low-radical-reuss}
\begin{aligned}
N&\ll_\varepsilon H^\varepsilon
 \min\left\{(\mathcal D\mathcal E M)^{1/2}+\mathcal D+\mathcal E,
              (UV M)^{1/2}+U+V\right\},\\
\log M&=\frac{9\log(\mathcal D\mathcal E)\log(UV)}{8\log x}.
\end{aligned}
\end{equation}
Write $\mathcal D\mathcal E=H^\theta$ up to constant factors. Since $x\asymp H$ and
$UV\asymp H^2/(\mathcal D\mathcal E)^2$, the determinant term in
\eqref{eq:low-radical-reuss} has exponent at most
\[
\phi(\theta)+O\left(\frac1{\log H}\right),\qquad
\phi(\theta)=\frac12\min(\theta,2-2\theta)+\frac98\theta(1-\theta),
\quad \frac7{12}\le\theta\le\frac34.
\]
The function increases up to $\theta=2/3$ and decreases afterwards, with
\[
\phi(2/3)=\frac13+\frac14=\frac7{12}.
\]

It remains to bound the additive terms. For $\theta\le2/3$, use the first
branch of \eqref{eq:low-radical-reuss}; by
\eqref{eq:low-radical-squareparts}, $\mathcal D+\mathcal E\ll H^{1/2}$. For $\theta\ge2/3$,
orient the cell so $\mathcal D\le \mathcal E$, hence $U\ge V$, and use the second branch.
If $U\le H^{7/12}$, its additive term is already small enough. Otherwise
$U=x/\mathcal D^2>H^{7/12}$ gives $\mathcal D\ll H^{5/24}$; together with
$\mathcal D\mathcal E\gg H^{2/3}$ this gives $\mathcal E\gg H^{11/24}$ and $V\ll H^{1/12}$.
Fixing $d,s$, the equation gives
\[
f^2s\equiv1\pmod{d^2},
\]
or $f^2s\equiv-1\pmod{d^2}$ in the reversed orientation. Since
$(d,s)=1$, each sign has at most $4^{\omega(d)}\ll_\varepsilon H^\varepsilon$
roots modulo $d^2$. Counting $f$ in its interval for each root and
summing over $d\ll \mathcal D$ and $s\ll V$ gives
\[
N\ll_\varepsilon H^\varepsilon
       \left(\frac{\mathcal E V}{\mathcal D}+\mathcal D V\right)
 \ll_\varepsilon H^\varepsilon
       \left(\frac{H}{\mathcal D\mathcal E}+\mathcal D V\right)
 \ll_\varepsilon H^{1/3+\varepsilon},
\]
since $\mathcal D V\ll H^{7/24}$. Thus all additive terms are below
$H^{7/12+\varepsilon}$. There are $O((\log H)^2)$ dyadic boxes and a
bounded number of cells in each. Summing
\eqref{eq:low-radical-pell-range} and the bounds above over all cells, the
logarithms are absorbed into $H^\varepsilon$ and
\begin{equation}\label{eq:low-radical-high-bound}
N_{\mathrm{hi}}\ll_\varepsilon H^{7/12+\varepsilon}.
\end{equation}
Combining \eqref{eq:low-radical-low-bound} and
\eqref{eq:low-radical-high-bound} proves the theorem.
\end{proof}

\begin{remark}
Diao's Lemma~2.1 \cite{Diao2026} gives the $7/12$ exponent for
$\delta_1x_1^2-\delta_2x_2^2=\vartheta$ in dyadic boxes with
$\Delta_1\le\Delta_2$, $\Delta_1\Delta_2\le B$, and
$\Delta_1X_1^2=\Delta_2X_2^2=B$, where $\vartheta\in\{\pm1,\pm2\}$. In our
cell the scales are $(\Delta_1,X_1,\Delta_2,X_2,B)
=(V,\mathcal E,U,\mathcal D,x)$, and the equation is
$sf^2-rd^2=1$. If the scales need to be ordered, interchanging the two
sides changes $\vartheta$ to $-1$. Since $\mathcal D\mathcal E\gg H^{7/12}$ and $x\asymp H$,
$UV=x^2/(\mathcal D\mathcal E)^2\ll H^{5/6}\ll x$ for sufficiently large $H$.
Thus Diao's condition $\Delta_1\Delta_2\le B$ holds. Diao proves Lemma~2.1 by
combining the determinant estimate (2.1) with the Fouvry--Jouve estimate
(2.2). The determinant input is the $(k,\ell)=(2,1)$ case of Reuss's
Theorem~1 \cite{Reuss2014}, applied above with $h=\pm1$. Diao cites this
estimate as Theorem~5 of Reuss's 2015 thesis; his lemma therefore gives the
high-product bound directly.
\end{remark}


\section*{Use of AI tools}

The author devised the main argument. Codex (OpenAI, GPT-6 Astra) carried
out the numerical checks of the identities and of the cubic family and
searched the literature. Claude Code (Anthropic, Claude Opus 5.5) checked
the arguments and the references. The author checked every step of the
proof and every reference, and takes full responsibility for the content
of the paper.

\section*{Statements and declarations}

\textbf{Competing interests.} The author has no competing interests.

\textbf{Funding.} No funding was received for this work.

\textbf{Data availability.} No data sets were generated or analyzed in this
work.

\begin{thebibliography}{10}

\bibitem{BSW1989}
R. Balasubramanian, T.~N. Shorey, and M. Waldschmidt,
\emph{On the maximal length of two sequences of consecutive integers with
the same prime divisors}, Acta Math. Hungar. \textbf{54} (1989), no.~3--4,
225--236, \url{https://doi.org/10.1007/BF01952052}.

\bibitem{Diao2026}
Y. Diao,
\emph{Class numbers and integer points on some Pellian surfaces},
J. Th\'eor. Nombres Bordeaux \textbf{37} (2025), no.~3, 973--988;
\url{https://doi.org/10.5802/jtnb.1348}; arXiv:2408.03774v2 (2026).

\bibitem{Erdos1980}
P. Erd\H{o}s,
\emph{How many pairs of products of consecutive integers have the same
prime factors?}, Amer. Math. Monthly \textbf{87} (1980), no.~5, 391--392,
\url{https://doi.org/10.1080/00029890.1980.11995045}.

\bibitem{ErdosStraus1977}
P. Erd\H{o}s and E.~G. Straus,
\emph{On products of consecutive integers}, in:
\emph{Number theory and algebra: collected papers dedicated to Henry B.
Mann, Arnold E. Ross, and Olga Taussky-Todd}, Academic Press, New York,
1977, 63--70.

\bibitem{Kane2011}
D.~M. Kane,
\emph{On a problem relating to the ABC conjecture},
preprint (2011), arXiv:1104.2635v1.

\bibitem{Kane2014}
D.~M. Kane,
\emph{On the number of ABC solutions with restricted radical sizes},
J. Number Theory \textbf{154} (2015), 32--43.
\url{https://doi.org/10.1016/j.jnt.2014.09.021}; preprint version:
arXiv:1104.2635v2 (2014). Proposition numbers refer to this version.

\bibitem{LL2025}
N. Lebowitz-Lockard,
\emph{On pairs of consecutive sequences with the same radicals},
preprint (2025), arXiv:2507.09899v1.

\bibitem{Lehmer1964}
D.~H. Lehmer,
\emph{On a problem of St{\o}rmer},
Illinois J. Math. \textbf{8} (1964), no.~1, 57--79,
\url{https://doi.org/10.1215/ijm/1256067456}.

\bibitem{Reuss2014}
T. Reuss,
\emph{Pairs of $k$-free numbers, consecutive square-full numbers},
preprint (2014), arXiv:1212.3150v2.

\bibitem{Sanna2022}
C. Sanna,
\emph{On the exponents in the factorizations of $r$ consecutive numbers},
Quaest. Math. \textbf{45} (2022), no.~8, 1221--1228.
\url{https://doi.org/10.2989/16073606.2021.1938277}.

\bibitem{Tao2026}
T. Tao,
\emph{Products of consecutive integers with unusual anatomy},
preprint (2026), arXiv:2603.27990v3.

\end{thebibliography}

\end{document}
